%% ma: L10-B12-0024
%% lop: 10
%% chuong: 5
%% bai: 12
%% dang: 10.12.6
%% loai: TN4
%% muc: TH
%% dap_an: "A"
%% nguon: "chuyen-de-cac-so-dac-trung-cua-mau-so-lieu-khong-ghep-nhom-toan-10-chuong-trinh-moi.pdf (GV. Phan Nhật Linh) (P01), Bài 01 Phần I Câu 4"
%% kien_thuc: "Lớp 10 Bài 12 (số gần đúng, sai số tuyệt đối, độ chính xác, sai số tương đối, quy tròn số)"
%% trang_thai: chinh-thuc
%% ghi_chu: "Sửa lời dẫn 'viết chuẩn giá trị gần đúng' thành 'số quy tròn' theo SGK KNTT; lời giải gốc không đúng quy tắc, đã viết lại. Tự duyệt theo yêu cầu giáo viên 22/10/2026."
\begin{ex}
Trong một thí nghiệm, hằng số $C$ được xác định là $5{,}73675$ với độ chính xác $d=0{,}00421$. Số quy tròn của $C$ là
\choice{\True $5{,}74$}{$5{,}736$}{$5{,}737$}{$5{,}7368$}
\loigiai{Chữ số khác $0$ đầu tiên của $d=0{,}00421$ ở hàng phần nghìn, nên quy tròn $C$ đến hàng phần trăm: $5{,}74$. Chọn A.}
\end{ex}
